% =============================================================================
% embedded-entwicklung.tex – Engineering-Regeln
% =============================================================================

\section{Embedded-Entwicklung – Engineering-Regeln}

\subsection{Begriffsrahmen und Zielsetzung}

Embedded-Entwicklung unterscheidet sich von Desktop-Software durch drei Eigenschaften: begrenzte Ressourcen, harte Echtzeitanforderungen, physische Schnittstellen. Die folgenden Regeln sind keine abstrakten Best Practices, sondern konkrete Antworten auf konkrete Probleme, die bei der Arbeit mit ESP32-S3, CD4021 und 74HC595 auftreten.

\begin{table}[htbp]
\centering
\begin{tabular}{@{}ll@{}}
\toprule
\textbf{Begriff} & \textbf{Definition} \\
\midrule
Determinismus & Gleiches Eingangssignal führt immer zum gleichen Verhalten \\
SSOT          & Single Source of Truth – jede Information hat genau eine Quelle \\
Hot-Path      & Code, der in jedem Scan-Zyklus ausgeführt wird \\
Edge          & Zustandsübergang (Flanke), nicht der Zustand selbst \\
Invariante    & Bedingung, die zu jedem Zeitpunkt gelten muss \\
\bottomrule
\end{tabular}
\caption{Begriffsrahmen Embedded-Entwicklung}
\label{tab:embedded-begriffe}
\end{table}

\textbf{Ziel:} Code, der deterministisch läuft, isoliert testbar ist und von 10 auf 100 Kanäle skaliert, ohne dass du die Logik ändern musst.

% -----------------------------------------------------------------------------
\subsection{Hardware ist die API}

\textbf{Problem:} In Desktop-Software definierst du Interfaces mit Methoden und Typen. In Embedded-Systemen ist die \glqq API\grqq{} physisch: Pins, Spannungspegel, Timing-Constraints. Diese Schnittstelle ist nicht verhandelbar.

\textbf{Regel:} Behandle Hardware-Konfiguration mit derselben Sorgfalt wie Software-Interfaces – dokumentiert, versioniert, an einer Stelle definiert.

\textbf{SSOT im Projekt:}

\begin{table}[htbp]
\centering
\begin{tabular}{@{}lll@{}}
\toprule
\textbf{Aspekt} & \textbf{SSOT-Datei} & \textbf{Beispielinhalt} \\
\midrule
Pins          & \code{firmware/include/hw/board\_pins.h} & \code{PIN\_SCK = D8} \\
Counts/Masks  & \code{firmware/include/hw/io\_counts.h}  & \code{LED\_COUNT = 10} \\
SPI-Parameter & \code{firmware/include/hw/spi\_params.h} & \code{CD4021\_FREQ = 500000} \\
\bottomrule
\end{tabular}
\caption{SSOT-Dateien für Hardware-Konfiguration}
\label{tab:ssot-dateien}
\end{table}

\textbf{Konsequenz:} Wenn du einen Pin änderst, änderst du ihn in \code{board\_pins.h} – nirgendwo sonst. Der Compiler erzwingt Konsistenz, weil alle anderen Dateien diese Header einbinden.

\textbf{Verboten:}

\begin{lstlisting}[style=arduino]
// Schlecht: Pin hardkodiert in Treiber
void cd4021_init() {
    pinMode(D1, OUTPUT);  // Wo kommt D1 her? Was ist D1?
}

// Gut: Pin aus SSOT
#include "hw/board_pins.h"
void cd4021_init() {
    pinMode(PIN_BTN_PS, OUTPUT);  // Selbstdokumentierend
}
\end{lstlisting}

% -----------------------------------------------------------------------------
\subsection{Determinismus vor Flexibilität}

\textbf{Problem:} Embedded-Systeme müssen in definierten Zeitfenstern reagieren. Ein Scan-Zyklus von \SI{10}{\milli\second} bedeutet: Nach \SI{10}{\milli\second} muss die nächste Messung starten – nicht \glqq ungefähr \SI{10}{\milli\second}\grqq, nicht \glqq wenn gerade Zeit ist\grqq.

\textbf{Regel:} Der Worst-Case zählt, nicht der Durchschnitt. Wenn eine Operation manchmal \SI{5}{\milli\second} und manchmal \SI{50}{\milli\second} dauert, ist sie im \SI{10}{\milli\second}-Zyklus unbrauchbar.

\textbf{Scan-Zyklus-Struktur:}

\begin{figure}[htbp]
\centering
\begin{tikzpicture}[
    phase/.style={draw, fill=CodeBg, minimum width=2.5cm, minimum height=1cm,
                  font=\small, rounded corners=2pt, align=center},
    arr/.style={-{Stealth[length=2mm]}, thick, color=ArduinoDark},
    label/.style={font=\footnotesize\ttfamily, color=MutedText}
]
\node[phase] (read) at (0,0) {read(raw)\\CD4021};
\node[phase] (update) at (3.5,0) {update(state)\\Debounce};
\node[phase] (derive) at (7,0) {derive(output)\\One-Hot};
\node[phase] (write) at (10.5,0) {write()\\HC595};

\draw[arr] (read) -- (update);
\draw[arr] (update) -- (derive);
\draw[arr] (derive) -- (write);

\node[label, below=0.2cm of read] {$\sim$\SI{20}{\micro\second}};
\node[label, below=0.2cm of update] {$\sim$\SI{5}{\micro\second}};
\node[label, below=0.2cm of derive] {$\sim$\SI{1}{\micro\second}};
\node[label, below=0.2cm of write] {$\sim$\SI{20}{\micro\second}};
\end{tikzpicture}
\caption{Scan-Zyklus-Struktur mit Timing}
\label{fig:scan-zyklus}
\end{figure}

\textbf{Timing-Budget (\SI{10}{\milli\second} = \SI{10000}{\micro\second}):}

\begin{table}[htbp]
\centering
\begin{tabular}{@{}lrr@{}}
\toprule
\textbf{Phase} & \textbf{Dauer} & \textbf{Anteil} \\
\midrule
read()      & $\sim$\SI{20}{\micro\second}   & 0,2\,\% \\
update()    & $\sim$\SI{5}{\micro\second}    & 0,05\,\% \\
derive()    & $\sim$\SI{1}{\micro\second}    & 0,01\,\% \\
write()     & $\sim$\SI{20}{\micro\second}   & 0,2\,\% \\
\textbf{Summe}   & $\sim$\SI{46}{\micro\second}   & $\sim$0,5\,\% \\
\textbf{Reserve} & $\sim$\SI{9954}{\micro\second} & 99,5\,\% \\
\bottomrule
\end{tabular}
\caption{Timing-Budget pro Scan-Zyklus}
\label{tab:timing-budget-scan}
\end{table}

\textbf{Was zerstört Determinismus:}

\begin{table}[htbp]
\centering
\begin{tabular}{@{}lll@{}}
\toprule
\textbf{Gut} & \textbf{Schlecht} & \textbf{Warum schlecht} \\
\midrule
\code{vTaskDelayUntil()} & \code{vTaskDelay()}        & Drift akkumuliert \\
Logging bei Events       & Logging pro Scan           & UART blockiert \\
Feste Zykluszeit         & \code{delay()} im Hot-Path & Timing wird unvorhersagbar \\
\bottomrule
\end{tabular}
\caption{Determinismus-Killer vermeiden}
\label{tab:determinismus}
\end{table}

\textbf{FreeRTOS-Beispiel:}

\begin{lstlisting}[style=arduino]
// Schlecht: Drift akkumuliert
void io_task(void* param) {
    while (true) {
        scan_cycle();
        vTaskDelay(pdMS_TO_TICKS(10));  // 10 ms NACH Ende des Zyklus
    }
}

// Gut: Feste Periode
void io_task(void* param) {
    TickType_t lastWake = xTaskGetTickCount();
    while (true) {
        scan_cycle();
        vTaskDelayUntil(&lastWake, pdMS_TO_TICKS(10));  // 10 ms AB Start
    }
}
\end{lstlisting}

% -----------------------------------------------------------------------------
\subsection{Edges statt Pegel}

\textbf{Problem:} Wenn du auf den Pegel \glqq Button ist gedrückt\grqq{} reagierst, feuert die Logik in jedem Scan-Zyklus. Bei \SI{10}{\milli\second} Intervall sind das 100 Events pro Sekunde, solange der Benutzer die Taste hält.

\textbf{Regel:} UI-Aktionen reagieren auf Flanken (Edges), nicht auf Pegel.

\textbf{Visualisierung:}

\begin{lstlisting}[style=wrap,numbers=none]
Pegel:    ****************************  (Taste 500 ms gehalten)
          ^                          ^
        Press                     Release

Edge:     ^                          ^   (zwei Events)
\end{lstlisting}

\textbf{Umsetzung im Projekt:}

\begin{center}
\code{CD4021 (Pegel) → Debounce (stabil) → Edge-Detection → Selection}
\end{center}

\textbf{Code-Beispiel:}

\begin{lstlisting}[style=arduino]
typedef struct {
    uint8_t stable;      // Letzter stabiler Zustand
    uint8_t previous;    // Fuer Edge-Detection
} ButtonState;

bool detectPressEdge(ButtonState* state, uint8_t current, uint8_t mask) {
    // Debounce hier ausgelassen fuer Klarheit

    bool wasPressed = (state->previous & mask) == 0;  // Active-low
    bool isPressed  = (current & mask) == 0;

    state->previous = current;

    // Press-Edge: war nicht gedrueckt, ist jetzt gedrueckt
    return (!wasPressed && isPressed);
}
\end{lstlisting}

\begin{tipbox}[Konsequenz]
Die Selection-Logik sieht nur \glqq Button 3 wurde gerade gedrückt\grqq, nicht \glqq Button 3 ist seit \SI{500}{\milli\second} gedrückt\grqq. Das macht das Verhalten vorhersagbar und testbar.
\end{tipbox}

% -----------------------------------------------------------------------------
\subsection{One-Hot ist Pflicht}

\textbf{Problem:} Bei einem Selection Panel mit 100 LEDs wäre es fatal, wenn versehentlich zwei LEDs gleichzeitig leuchten. Das kann passieren, wenn du LEDs an verschiedenen Stellen im Code setzt/löschst.

\textbf{Regel:} One-Hot-Zustände entstehen an genau einer Stelle im Code – nicht verteilt.

\textbf{Invariante:}

\begin{itemize}
\item \code{selectedId == 0} → alle LEDs aus
\item \code{selectedId} $\in$ \code{[1..N]} → genau LED \code{selectedId} an
\end{itemize}

\textbf{Zentrale Implementierung:}

\begin{lstlisting}[style=arduino]
void setOneHot(uint8_t selectedId, uint8_t* ledState, size_t bytes) {
    // 1. ALLE Bits loeschen (garantiert sauberen Zustand)
    memset(ledState, 0x00, bytes);

    // 2. Genau EIN Bit setzen (wenn ID > 0)
    if (selectedId >= 1 && selectedId <= LED_COUNT) {
        uint8_t byteIdx = (selectedId - 1) / 8;
        uint8_t bitPos  = 7 - ((selectedId - 1) % 8);  // MSB-first
        ledState[byteIdx] |= (1u << bitPos);
    }

    // 3. Ungueltige Bits im letzten Byte maskieren
    ledState[bytes - 1] &= LED_LAST_MASK;
}
\end{lstlisting}

\textbf{Warum} \code{memset} \textbf{zuerst?}

Alternative Ansätze wie \glqq altes Bit löschen, neues Bit setzen\grqq{} können bei Bugs, Race Conditions oder vergessenen Aufrufen zu Zuständen mit zwei gesetzten Bits führen:

\begin{lstlisting}[style=arduino]
// Gefaehrlich: Zwei Operationen, dazwischen kann etwas schiefgehen
ledState[oldByte] &= ~oldMask;  // Was wenn hier ein Interrupt kommt?
ledState[newByte] |= newMask;   // Jetzt sind kurzzeitig zwei Bits gesetzt

// Sicher: Atomare Operation auf dem Zustand
memset(ledState, 0, bytes);     // Alles aus
ledState[byteIdx] |= mask;      // Genau eins an
\end{lstlisting}

\begin{warnbox}[Verboten]
\begin{lstlisting}[style=arduino,numbers=none]
// Schlecht: LEDs "nebenbei" setzen
void handleSpecialCase() {
    ledState[0] |= 0x80;  // LED 1 zusaetzlich an - bricht One-Hot!
}
\end{lstlisting}
\end{warnbox}

% -----------------------------------------------------------------------------
\subsection{Skalierung durch Formeln}

\textbf{Problem:} Code mit hardkodierten Werten wie \code{uint8\_t leds[2]} oder \code{for (int i = 0; i < 10; i++)} muss bei Skalierung manuell angepasst werden. Das ist fehleranfällig.

\textbf{Regel:} Alle Größen, die von der Kanalzahl abhängen, werden aus \code{*\_COUNT} abgeleitet.

\textbf{SSOT-Definitionen:}

\begin{lstlisting}[style=arduino]
// io_counts.h
#define BTN_COUNT  10
#define LED_COUNT  10

// Abgeleitete Werte (automatisch korrekt)
#define BTN_BYTES  ((BTN_COUNT + 7) / 8)  // = 2
#define LED_BYTES  ((LED_COUNT + 7) / 8)  // = 2

// Maske fuer ungueltige Bits im letzten Byte
// Bei 10 Kanaelen: 10 % 8 = 2 gueltige Bits -> Maske = 0xC0
#define LED_LAST_MASK  (0xFF << (8 - (LED_COUNT % 8 == 0 ? 8 : LED_COUNT % 8)))
\end{lstlisting}

\textbf{Vergleich:}

\begin{table}[htbp]
\centering
\begin{tabular}{@{}lll@{}}
\toprule
\textbf{Richtig} & \textbf{Falsch} & \textbf{Warum falsch} \\
\midrule
\code{LED\_COUNT}     & \code{10}   & Bei Änderung vergessen \\
\code{LED\_BYTES}     & \code{2}    & Inkonsistent bei Skalierung \\
\code{LED\_LAST\_MASK} & \code{0xC0} & Muss neu berechnet werden \\
\bottomrule
\end{tabular}
\caption{Richtige vs. falsche Parametrisierung}
\label{tab:parametrisierung}
\end{table}

\textbf{Skalierungsbeispiel:}

\begin{table}[htbp]
\centering
\begin{tabular}{@{}lrll@{}}
\toprule
\textbf{Parameter} & \textbf{10 Kanäle} & \textbf{100 Kanäle} & \textbf{Änderung nötig} \\
\midrule
\code{LED\_COUNT}     & 10        & 100       & ✓ (einmal) \\
\code{LED\_BYTES}     & 2         & 13        & Automatisch \\
\code{LED\_LAST\_MASK} & \code{0xC0} & \code{0xF0} & Automatisch \\
Logik-Code            & –         & –         & Keine \\
\bottomrule
\end{tabular}
\caption{Skalierung von 10 auf 100 Kanäle}
\label{tab:skalierung-beispiel}
\end{table}

\begin{tipbox}[Konsequenz]
Skalierung von 10 auf 100 Kanäle erfordert eine Zeile Änderung in \code{io\_counts.h}.
\end{tipbox}

% -----------------------------------------------------------------------------
\subsection{Schichten einhalten}

\textbf{Problem:} Wenn Treiber Geschäftslogik enthalten und Logik Hardware-Zugriffe macht, wird das System untestbar und unwartbar.

\textbf{Regel:} Abhängigkeiten zeigen nach unten. Höhere Schichten nutzen niedrigere, nie umgekehrt.

\begin{figure}[htbp]
\centering
\begin{tikzpicture}[
    layer/.style={draw, fill=CodeBg, minimum width=11cm, minimum height=1cm,
                  font=\small, rounded corners=2pt, align=center},
    arr/.style={-{Stealth[length=2mm]}, thick, color=ArduinoDark}
]
\node[layer] (app) at (0,3) {\textbf{app/}\hspace{1cm}Orchestrierung (Loop/Task, Timing) – Ruft auf, implementiert nichts};
\node[layer] (logic) at (0,2) {\textbf{logic/}\hspace{1cm}Regeln (Debounce, Selection, Pipeline) – Keine Hardware-Abhängigkeiten};
\node[layer] (drivers) at (0,1) {\textbf{drivers/}\hspace{1cm}IC-Zugriff (CD4021, HC595) – Keine Geschäftslogik};
\node[layer] (hal) at (0,0) {\textbf{hal/}\hspace{1cm}Hardware-Abstraktion (SPI-Transfer) – Plattform-spezifisch};
\draw[arr] (app.south) -- (logic.north);
\draw[arr] (logic.south) -- (drivers.north);
\draw[arr] (drivers.south) -- (hal.north);
\end{tikzpicture}
\caption{Schichtenmodell mit Verantwortlichkeiten}
\label{fig:schichten-embedded}
\end{figure}

\textbf{Was darf wo?}

\begin{table}[htbp]
\centering
\begin{tabular}{@{}lll@{}}
\toprule
\textbf{Schicht} & \textbf{Darf} & \textbf{Darf nicht} \\
\midrule
\code{app/}     & Schichten orchestrieren, Timing      & Logik implementieren \\
\code{logic/}   & Zustände verarbeiten, Entscheidungen & Pins kennen, SPI aufrufen \\
\code{drivers/} & Bytes lesen/schreiben                & Debounce-Entscheidungen \\
\code{hal/}     & Register beschreiben                 & IC-Protokolle kennen \\
\bottomrule
\end{tabular}
\caption{Zuständigkeiten der Schichten}
\label{tab:schichten-embedded}
\end{table}

\textbf{Warum ist das wichtig?}

Die \code{logic/}-Schicht ist hardware-unabhängig. Du kannst sie auf dem Host-PC testen:

\begin{lstlisting}[style=arduino]
// test_selection.cpp - laeuft ohne ESP32
void test_selection_toggles_on_repress() {
    SelectionState state = {0};

    // Simuliere: Button 3 gedrueckt
    selection_update(&state, 3, true);
    TEST_ASSERT_EQUAL(3, state.selectedId);

    // Simuliere: Button 3 nochmal gedrueckt -> Toggle aus
    selection_update(&state, 3, true);
    TEST_ASSERT_EQUAL(0, state.selectedId);
}
\end{lstlisting}

% -----------------------------------------------------------------------------
\subsection{Testen ist Pflicht}

\textbf{Problem:} Embedded-Bugs sind schwer zu debuggen. Du siehst nicht, was auf der Leitung passiert. Ohne Tests merkst du Regressionen erst, wenn das Gerät falsch reagiert.

\textbf{Regel:} Jede Änderung erfordert den passenden Test.

\textbf{Zweistufige Strategie:}

\begin{table}[htbp]
\centering
\begin{tabular}{@{}lll@{}}
\toprule
\textbf{Test-Typ} & \textbf{Befehl} & \textbf{Testet} \\
\midrule
Unit & \code{pio test -e seeed\_xiao\_esp32s3\_test -v} & Logik isoliert \\
E2E  & \code{pio test -e seeed\_xiao\_esp32s3\_e2e -f test\_e2e\_hw -v} & Hardware-Integration \\
\bottomrule
\end{tabular}
\caption{Test-Strategie}
\label{tab:test-strategie}
\end{table}

\textbf{Wann welcher Test?}

\begin{table}[htbp]
\centering
\begin{tabular}{@{}lcc@{}}
\toprule
\textbf{Änderung} & \textbf{Unit-Test} & \textbf{E2E-Test} \\
\midrule
Debounce-Algorithmus & ✓ & – \\
Selection-Policy     & ✓ & – \\
Pin-Zuordnung        & – & ✓ \\
SPI-Timing           & – & ✓ \\
Neuer Kanal-Count    & ✓ & ✓ \\
\bottomrule
\end{tabular}
\caption{Test-Zuordnung nach Änderungstyp}
\label{tab:test-zuordnung}
\end{table}

\textbf{Unit-Test-Beispiel:}

\begin{lstlisting}[style=arduino]
void test_debounce_filters_noise() {
    DebounceState state = {.stable = 0xFF, .counter = 0};

    // Prellen simulieren: 1 -> 0 -> 1 -> 0 -> 0 -> 0
    TEST_ASSERT_FALSE(debounce_update(&state, 0xFE, 3));  // Zaehler = 1
    TEST_ASSERT_FALSE(debounce_update(&state, 0xFF, 3));  // Reset
    TEST_ASSERT_FALSE(debounce_update(&state, 0xFE, 3));  // Zaehler = 1
    TEST_ASSERT_FALSE(debounce_update(&state, 0xFE, 3));  // Zaehler = 2
    TEST_ASSERT_TRUE(debounce_update(&state, 0xFE, 3));   // Zaehler = 3 -> Edge!
}
\end{lstlisting}

% -----------------------------------------------------------------------------
\subsection{Debugging-Strategie}

\textbf{Regel:} Debugge von unten nach oben – Hardware vor Software.

\textbf{Reihenfolge:}

\begin{enumerate}
\item \textbf{Verkabelung} – Spannung, GND, Signale korrekt?
\item \textbf{SPI-Signale} – Clock, MOSI, MISO sichtbar?
\item \textbf{Raw-Daten} – Erwartete Bytes vom CD4021?
\item \textbf{Logik-Zustand} – Debounce, Selection korrekt?
\item \textbf{Output} – LEDs wie erwartet?
\end{enumerate}

\textbf{Werkzeuge:}

\begin{table}[htbp]
\centering
\begin{tabular}{@{}lll@{}}
\toprule
\textbf{Ebene} & \textbf{Werkzeug} & \textbf{Prüfung} \\
\midrule
Verkabelung & Multimeter     & VCC = \SI{3,3}{\volt}? GND verbunden? \\
SPI         & Logic Analyzer & Timing, Pegel, Protokoll \\
Raw-Daten   & Serial Monitor & \code{RAW: 0xFD 0xFF} \\
Logik       & Unit-Tests     & Algorithmen isoliert \\
Output      & Augen          & LED 5 leuchtet? \\
\bottomrule
\end{tabular}
\caption{Debugging-Werkzeuge nach Ebene}
\label{tab:debug-werkzeuge}
\end{table}

\textbf{Typische Fehlerbilder:}

\begin{table}[htbp]
\centering
\begin{tabular}{@{}lll@{}}
\toprule
\textbf{Symptom} & \textbf{Wahrscheinliche Ursache} & \textbf{Prüfung} \\
\midrule
RAW = 0xFF (immer) & Pullups fehlen oder PS falsch & Verkabelung \\
RAW = 0x00 (immer) & Kurzschluss nach GND          & Verkabelung \\
Falsche ID erkannt & MSB/LSB vertauscht            & Bit-Mapping prüfen \\
LEDs flackern      & Timing zu schnell             & SPI-Frequenz senken \\
Doppelte Events    & Debounce zu kurz              & Threshold erhöhen \\
\bottomrule
\end{tabular}
\caption{Typische Fehlerbilder und Diagnose}
\label{tab:fehlerbilder}
\end{table}

% -----------------------------------------------------------------------------
\subsection{Häufige Fehler}

\begin{table}[htbp]
\centering
\begin{tabularx}{\textwidth}{@{}lXX@{}}
\toprule
\textbf{Fehler} & \textbf{Konsequenz} & \textbf{Lösung} \\
\midrule
\code{delay()} im Scan    & Timing gestört & \code{millis()}-basiert oder \code{vTaskDelayUntil()} \\
Logging pro Scan          & Serial überlastet (\SI{4}{\milli\second} pro Zeile) & Nur bei Events loggen \\
Globale Variablen         & Race Conditions bei Interrupts & Zentraler State mit Mutex \\
Feste Array-Größen        & Nicht skalierbar & \code{*\_BYTES} Makros verwenden \\
LEDs verteilt setzen      & One-Hot gebrochen & Zentrale \code{setOneHot()} Funktion \\
Pin in Treiber            & SSOT verletzt & \code{board\_pins.h} verwenden \\
\bottomrule
\end{tabularx}
\caption{Häufige Fehler und Lösungen}
\label{tab:haeufige-fehler}
\end{table}

% -----------------------------------------------------------------------------
\subsection{Checkliste vor Commit}

\begin{itemize}
\item[$\square$] Pins/Counts nur in SSOT-Dateien (\code{include/hw/*})
\item[$\square$] Kein \code{delay()} im Scan-Pfad
\item[$\square$] One-Hot zentral in einer Funktion erzeugt
\item[$\square$] Letztes Byte mit \code{*\_LAST\_MASK} maskiert
\item[$\square$] \code{logic/} enthält keine Hardware-Includes
\item[$\square$] Unit-Tests grün: \code{pio test -e seeed\_xiao\_esp32s3\_test -v}
\item[$\square$] Bei HW-Änderung: E2E grün: \code{pio test -e seeed\_xiao\_esp32s3\_e2e -f test\_e2e\_hw -v}
\item[$\square$] Dokumentation aktualisiert (bei Änderungen an Wiring, Architektur)
\end{itemize}

% -----------------------------------------------------------------------------
\subsection{Weiterführende Dokumentation}

\begin{description}
\item[\code{firmware/docs/architecture.md}] Schichtenmodell-Details
\item[\code{firmware/docs/wiring.md}] Hardware-Verdrahtung
\item[\code{firmware/docs/test-plan.md}] Testabdeckung
\item[\code{firmware/docs/debug-playbook.md}] Erweiterte Fehlerbehebung
\item[\code{docs/embedded-grundlagen.md}] Konzeptuelle Grundlagen
\end{description}
